\documentclass[UTF8,fontset=macnew,11pt]{ctexart}

\usepackage[a4paper,margin=2.25cm]{geometry}
\usepackage{amsmath,amssymb}
\usepackage{booktabs}
\usepackage{tabularx}
\usepackage{array}
\usepackage{float}
\usepackage{xcolor}
\usepackage{hyperref}

\hypersetup{
  colorlinks=true,
  linkcolor=blue!50!black,
  citecolor=blue!50!black,
  urlcolor=blue!50!black
}

\setlength{\parskip}{0.36em}
\setlength{\parindent}{2em}
\setlength{\emergencystretch}{2em}

\newcommand{\method}{Resource-NMCTS}
\newcommand{\pareto}{Pareto-Resource-NMCTS}
\newcommand{\screen}{Boolean-ring screen}
\newcommand{\score}{\ensuremath{\mathrm{score}}}
\newcommand{\anf}{\ensuremath{\mathrm{ANF}}}
\newcommand{\fprm}{\ensuremath{\mathrm{FPRM}}}
\newcommand{\mcts}{\ensuremath{\mathrm{MCTS}}}
\newcommand{\xor}{\ensuremath{\oplus}}
\newcommand{\ttable}{truth-table}

\title{面向资源约束量子布尔函数综合的神经蒙特卡洛树搜索方法}
\author{中文论文稿 v11：主线收束版}
\date{2026年7月8日}

\begin{document}
\pagestyle{plain}
\maketitle

\begin{abstract}
量子布尔函数 oracle 综合是 Grover 搜索、密码分析和若干量子算法子程序中的基础编译步骤。容错量子计算中 T 门和多控门分解代价较高，因而 oracle 的 T-count、CNOT、深度和辅助比特会直接影响逻辑层资源估计。本文将布尔 oracle 综合表述为代数正规形（algebraic normal form, \anf{}）和固定极性 Reed--Muller（fixed-polarity Reed--Muller, \fprm{}）项集合上的资源加权搜索问题，提出结合神经动作先验、蒙特卡洛树搜索、布尔环线性因子筛选、结构级 depth policy 与保守 guard 的 \method{} 框架。该方法不依赖 SSHR 的 parallelotope cover 空间；SSHR 仅作为 CNOT-oriented small-scale baseline 使用。实验表明，在 $n\leq6$ 的 177 个传统函数上，\pareto{} 相对 ESOP cube beam 获得 174/0/3 的 score 胜/负/平，平均 score 降低 36.09\%；相对 weighted ESOP-MILP 获得 167/3/7，平均 score 降低 29.84\%。在 $n=18$ 的 12 个随机 \anf{} 函数上，adaptive depth-2 \screen{} 相对 single screen 平均 score 降低 6.47\%，完整 \method{} 相对 adaptive screen 进一步降低 23.85\%。在 $n=20$ 边界函数上，screen-gated \method{} 与完整 \method{} 资源完全持平，并把平均运行时间从 77.10 s 降至 20.71 s；在独立 $n=19,20$ gate holdout 中，screen-gated \method{} 在 16 个函数上全部 score 持平并平均节省 36.83\% 时间。在 $n=20,22,24,28$ 的 192 个项集级 scale 测试中，depth policy 相对 single screen 获得 169/0/23，平均 score 降低 6.56\%；相对 all-depth adaptive 只损失 0.08\% 平均 score，并节省 31.01\% 平均时间。进一步在 $n=32,36,40$ 的 144 个项集上，depth policy 相对 single screen 获得 110/0/34、平均 score 降低 5.55\%，并与 all-depth adaptive 全部 score 持平、平均节省 33.14\% 时间。高预算 depth-frontier 实验显示，depth-4 \screen{} 相对 fixed depth-2 在 $n=20,28,40$ 上获得 49/0/23，平均 score 进一步降低 3.10\%，但平均运行时间增加 682.97\%。上述项集级实验中，3000/3000 个方法行均通过 factor plan 的 \anf{} 展开验证和 emitted X/CNOT/MCT circuit 的 GF(2) 多项式符号模拟验证。结果支持一条独立于 SSHR 的资源约束 oracle 综合主线，同时也明确边界：本文只讨论逻辑层综合，不包含硬件 mapping；$n>20$ 结果是项集级符号验证，不是完整 \ttable{} simulation；当前高维 learned prior 的独立质量收益仍弱，结构级 AI 仍需从安全门控推进到更大幅度的质量收益。
\end{abstract}

\noindent\textbf{关键词：}量子 oracle 综合；布尔函数综合；神经蒙特卡洛树搜索；ANF；FPRM；布尔环；资源约束综合

\section{引言}

给定布尔函数 $f:\{0,1\}^n\rightarrow\{0,1\}$，量子 oracle 通常实现为
\begin{equation}
  |x\rangle |y\rangle \mapsto |x\rangle |y\xor f(x)\rangle .
  \label{eq:oracle}
\end{equation}
这一变换是许多量子算法调用经典谓词的入口。若 $f$ 被直接展开为单项式并逐项写入输出位，线路实现简单且容易验证，但会忽略项之间的可共享结构；若过度追求共享，又可能增加辅助比特、CNOT 和调度深度。因此，量子布尔函数综合不适合只写成 ``最小门数'' 问题，而应作为多资源约束下的组合搜索问题处理。

本文的出发点是 \anf{} 表示。若
\begin{equation}
  f(x)=\bigoplus_{m\in\mathcal{M}}\prod_{i\in m}x_i ,
  \label{eq:anf}
\end{equation}
则每个单项式都可转化为受控乘法结构，并异或到输出位。\anf{} 的优势是语义透明；弱点是它不自动暴露公共因子、线性奇偶量、极性重写和递归子问题。已有 XAG、ROS、BDD 和 SSHR 等路线分别从 XOR/AND 图、资源约束 LUT 映射、决策图和 parallelotope cover 改进 oracle 或可逆逻辑综合\cite{meuli2022xag,meuli2020ros,wille2009bdd,zheng2025sshr}。这些工作说明中间表示会显著改变可优化资源，但仍留下一个问题：能否直接在 \anf{}/\fprm{} 项集合上搜索可复用代数结构，并用 AI 策略在该搜索空间中选择更优动作？

本文的一句话论证是：在逻辑层量子布尔函数 oracle 综合中，\anf{}/\fprm{} 项集合上的资源感知神经搜索，配合布尔环线性因子筛选和保守结构门控，可以在传统函数与高维随机项集合上降低 T-count 和加权逻辑资源；当前边界是该方法尚未覆盖硬件 mapping、真实噪声模型和完整的大规模 \ttable{} 语义验证。本文不把目标写成 ``AI-SSHR''，因为方法本身并不继承 SSHR 的 parallelotope 空间。SSHR 在本文中只是小规模 CNOT 导向基线，用于暴露 T-count 与 CNOT 指标之间的取舍。

\section{相关工作与本文定位}

\subsection{布尔表示与 oracle 综合}

Xor-And-Inverter Graphs（XAG）把 XOR 与 AND 结构分离，使 AND 节点数量与非 Clifford 成本直接相关。Meuli 等人将 XAG 用于 quantum compilation\cite{meuli2022xag}，并从 multiplicative complexity 角度解释低 T-count oracle 的结构来源\cite{meuli2019multiplicative}。ROS 进一步把 oracle synthesis 建模为 resource-constrained LUT mapping\cite{meuli2020ros}。这些工作与本文的共同点是把布尔结构视为资源优化入口；不同点是，本文不固定使用 XAG 或 LUT，而是在 \anf{}/\fprm{} 项集合上搜索因子分解与 compute/uncompute 结构。

BDD-based reversible synthesis 用决策图处理较大布尔函数\cite{wille2009bdd}。Back-end-aware oracle synthesis 则把 XAG 优化与后端质量指标连接起来\cite{yu2025backend}。本文与这些工作互补：当前实验只报告逻辑层 T-count、CNOT、depth、gate count 和辅助比特，不报告设备连通性、routing、magic-state factory 或噪声下保真度。

SSHR 使用 parallelotope 的空间结构做 small-scale Boolean function 的 CNOT-oriented synthesis\cite{zheng2025sshr}。本文复现并保留 SSHR-H/SSHR-I 作为小规模比较对象，但方法本体不枚举 parallelotope candidate，也不以 CNOT-only 最优为唯一目标。更准确的定位是：本文是一条面向低 T-count 和加权逻辑资源的 \anf{}/\fprm{} 搜索路线，SSHR 是其中一个能够说明 CNOT trade-off 的外部参照。

\subsection{学习式搜索与量子线路优化}

学习式搜索已经用于若干离散综合问题。ShortCircuit 使用 AlphaZero 风格搜索设计经典布尔电路\cite{tsaras2024shortcircuit}；Gumbel AlphaZero 被用于 Clifford+T unitary synthesis\cite{rietsch2024unitary}；ZX-calculus 优化中出现了强化学习和图神经网络驱动的 rewrite-rule 选择\cite{riu2025rlzx}；MonteQ 使用 \mcts{} 处理量子线路综合中的动作排序\cite{machiya2026monteq}。这些工作说明学习式策略适合组合空间，但它们的状态和动作并不是 Boolean oracle 的 \anf{}/\fprm{} 因子分解。本文的 AI 贡献应被理解为面向 Boolean oracle 结构的动作排序、结构深度选择和安全门控，而不是泛化的 ``RL 直接生成量子线路''。

\section{问题定义与资源模型}

本文输入为布尔函数的 \ttable{} 或 \anf{} 表示，输出为实现式~\eqref{eq:oracle} 的逻辑层可逆线路。候选线路由 X、CNOT 和多控 Toffoli（MCT）抽象门组成，并在资源统计阶段转化为 T-count、CNOT、深度、总 gate count 和 peak ancilla。用于搜索排序的统一分数为
\begin{equation}
  \score = T + 0.04\,\mathrm{CNOT} + 0.015\,\mathrm{depth}
         + 0.01\,\mathrm{gates} + 2\,\mathrm{ancilla}.
  \label{eq:score}
\end{equation}
该分数只用于逻辑层候选选择，不代表某个具体硬件平台的物理开销。为避免单一分数掩盖取舍，实验表同时报告各项资源。

在 $n\leq20$ 的函数级实验中，本文对候选线路做完整 \ttable{} 语义验证。在 $n>20$ 的项集级 scale 实验中，完整 \ttable{} 构造本身不可接受，因此采用两层符号验证：第一层把 factor plan 展开回 \anf{} 项集合，检查目标项集合等价；第二层对 emitted X/CNOT/MCT circuit 做 GF(2) 多项式符号模拟，检查输入线复原、输出线等价和辅助线清零。该验证不等同于逐点枚举 $2^n$ 个输入，但比只统计 plan 资源更强，因为它验证了生成线路的代数语义。

\begin{table}[t]
\centering
\small
\caption{本文术语与边界。}
\label{tab:terms}
\begin{tabularx}{\linewidth}{>{\raggedright\arraybackslash}p{0.25\linewidth}>{\raggedright\arraybackslash}X}
\toprule
术语 & 本文含义 \\
\midrule
\method{} & 在 \anf{}/\fprm{} 项集合上进行资源加权搜索的主框架，包括神经先验、\mcts{}、guard 和 portfolio 选择。 \\
\screen{} & 在 square-free Boolean ring 中搜索可复用线性因子的结构筛选分支，允许线性因子与 quotient 共享变量。 \\
Depth policy & 预测 single、depth-1 或 depth-2 screen 的结构级神经策略。 \\
Skip guard & 判断是否安全跳过更深 screen 或 Resource tail 的保守门控；当前目标优先是 zero-loss。 \\
SSHR & CNOT-oriented small-scale baseline；本文方法不从 SSHR 搜索空间出发。 \\
\bottomrule
\end{tabularx}
\end{table}

\section{方法}

\subsection{状态、动作与资源感知搜索}

\method{} 的状态是当前尚未综合的项集合 $\mathcal{M}$。一次动作选择一个可计算、可反算的局部代数结构，将 $\mathcal{M}$ 分解为 quotient 子问题和 rest 子问题，再递归生成 compute/action/uncompute 形式的线路。候选动作包括公共单项式因子、\fprm{} 极性重写后的公共项、线性奇偶因子、布尔环线性因子，以及由神经模型扩展的 root actions。

神经动作先验并不直接输出量子门序列，而是对候选因子动作排序。特征包括候选覆盖率、项数变化、项平均次数、变量覆盖密度、即时资源收益和子问题规模等结构统计。\mcts{} 在有限预算内探索候选组合，rollout 使用确定性资源估计与已有 baseline 候选。最终选择由 baseline-preserving guard 兜底：若学习候选不优于确定性 baseline，则返回 baseline。这个设计牺牲部分探索自由度，但能显著降低 AI 负迁移造成资源退化的风险。

\subsection{布尔环线性因子}

早期 linear-pair 分支主要搜索形如
\begin{equation}
  (x_i\xor x_j\xor\cdots)\cdot h(x)
  \label{eq:linear-factor}
\end{equation}
的结构，并要求线性因子变量与 quotient $h(x)$ 的变量不相交。该约束便于实现，却排除了 square-free Boolean ring 中合法的重叠变量情形。本文的 \screen{} 允许二者共享变量，并在展开时使用 $x_i^2=x_i$ 与 GF(2) 偶数项消去。例如
\begin{equation}
  (x_i\xor x_j)x_i m = x_i m \xor x_i x_j m .
  \label{eq:boolean-ring}
\end{equation}
线路层面仍可先用 CNOT 计算线性奇偶量，再用该辅助量控制 quotient 子线路。因此，布尔环线性因子不是简单扩大 beam 宽度，而是把原先被实现约束排除的代数候选重新纳入可综合搜索空间。

\subsection{结构级 AI 策略}

高维实验显示，fixed single screen、depth-1 screen 和 depth-2 screen 的收益依赖项集合结构。本文因此训练 depth policy，输入项数、次数分布、变量覆盖密度、二元共现密度和 root Boolean-ring action 统计，输出应该使用的 screen 深度。监督标签来自实际运行三种 screen 后按式~\eqref{eq:score} 选择的最佳 depth。

由于固定 depth-2 screen 是很强的质量 baseline，本文又训练保守 depth-2 skip guard。Static-direct 模式只使用项集合静态特征，若预测 depth-1 足以与 fixed depth-2 持平，则直接运行较浅 screen；否则运行 depth-2。Shallow-staged 模式先运行 single/depth-1 screen，再用浅层观测判断是否跳过 depth-2。阈值按 zero-false-skip 约束选择，使门控优先避免质量损失。

此外，$n=20$ 边界切片中完整 \method{} 的 Resource tail 与 adaptive screen 资源持平但运行更慢。为控制该长尾，本文训练 screen-gate 判断是否可以跳过 Resource tail。该门控当前只作为运行时控制证据：它说明某些切片上能以零资源损失节省时间，但还不能说明高维 Resource tail 普遍不必要。

\section{实验设置}

实验分为五组。第一组使用 $n\leq6$ 的 177 个传统函数，与 direct ANF、AND-direct ANF、ESOP cube beam、weighted ESOP-MILP、SSHR-H、ABC 和 BDD baseline 比较。第二组使用 $n=18$ 随机 \anf{} 函数，观察 adaptive screen 与完整 \method{} 的差距。第三组使用 $n=20$ 边界函数，检验在 FPRM root beam 和 fast linear-pair 超时区域中，depth-2 screen 是否仍能给出可运行改进，并用 $n=19,20$ held-out 切片验证 screen-gate 是否会错误跳过 Resource tail。第四组使用 $n=14,16,18$ 生成式项集训练 depth policy 与 skip guard，并在 held-out $n=20$ 项集上测试结构选择能力。第五组在 \anf{} 项集合上评估 $n=20,22,24,28$ 和扩展 $n=32,36,40$ 的 screen-scale 行为，并对每个 plan 和 emitted circuit 做符号等价验证。

所有 paired comparison 只纳入函数名或项集名匹配、语义正确且未 timeout 的结果。Timeout 行保留为运行边界证据，但不参与正确结果均值比较。本文所有实验均为逻辑层资源估计，不包含 hardware mapping 和 noise-aware scheduling。

\section{结果}

\subsection{小规模传统函数}

表~\ref{tab:traditional} 给出 $n\leq6$ 传统函数上的平均逻辑资源。相对 direct ANF，\pareto{} 平均 T-count 降低 72.25\%，平均 score 降低 67.80\%。相对 ESOP cube beam，\pareto{} 获得 174/0/3 的 score 胜/负/平，平均 score 降低 36.09\%。相对 weighted ESOP-MILP，获得 167/3/7，平均 score 降低 29.84\%。

\begin{table}[t]
\centering
\small
\caption{$n\leq6$ 传统函数切片上的平均逻辑资源。}
\label{tab:traditional}
\begin{tabular}{lrrrrr}
\toprule
方法 & 平均 T & 平均 CNOT & 平均 depth & 平均 ancilla & 平均 score \\
\midrule
Direct ANF & 184.1 & 134.2 & 134.6 & 1.33 & 194.3 \\
AND-direct ANF & 101.0 & 166.5 & 166.9 & 2.18 & 115.2 \\
\method{} & 43.9 & 89.9 & 94.5 & 1.92 & 53.2 \\
\pareto{} & \textbf{40.4} & 83.0 & \textbf{88.0} & 2.03 & \textbf{49.6} \\
ESOP-MILP & 83.6 & 133.5 & 159.6 & 2.32 & 96.7 \\
SSHR-H & 81.0 & \textbf{67.6} & 88.7 & 1.36 & 88.2 \\
\bottomrule
\end{tabular}
\end{table}

该结果支持本文的低 T-count 与加权资源优势，但也给出重要边界：SSHR-H 的平均 CNOT 最低，因此本文不能主张 CNOT-only 指标全面支配 SSHR。更准确的表述是，本文方法在当前资源权重下用一定 CNOT 与辅助比特 trade-off 换取更低 T-count 和更低 score。

\subsection{高维函数级边界}

表~\ref{tab:n18n20} 汇总 $n=18$ 与 $n=20$ 函数级边界结果。在 $n=18$ 随机 \anf{} 函数上，adaptive depth-2 \screen{} 相对 single screen 获得 11/0/1 的 score 胜/负/平，平均 score 从 11349.06 降至 10670.94；完整 \method{} 平均 score 为 7315.62，相对 adaptive screen 进一步降低 23.85\%。这说明 $n=18$ 上快速 screen 是有效候选生成器，但不能替代完整资源搜索。

在 $n=20$ 边界函数上，FPRM root beam 和 fast linear-pair 均触发 300 s task timeout。Adaptive depth-2 screen 相对 single screen 获得 5/0/1，平均 score 降低 7.47\%。集成该分支后，\method{} 相对 AND-direct ANF 获得 5/0/1，平均 score 降低 11.80\%。由于此切片中完整 Resource tail 与 adaptive screen 资源持平，screen-gated \method{} 可以把平均时间从 77.10 s 降到 20.71 s。

\begin{table}[t]
\centering
\scriptsize
\caption{高维函数级边界结果。}
\label{tab:n18n20}
\begin{tabular}{llrrrr}
\toprule
切片 & 方法 & 平均 T & 平均 CNOT & 平均 score & 平均时间(s) \\
\midrule
$n=18$ & Single screen & 10389.33 & 16268.42 & 11349.06 & 0.82 \\
$n=18$ & Adaptive depth-2 screen & 9767.00 & 15301.58 & 10670.94 & 4.39 \\
$n=18$ & \method{} & \textbf{6639.33} & \textbf{11380.92} & \textbf{7315.62} & 226.18 \\
$n=20$ & AND-direct ANF & 21516.67 & 33635.67 & 23491.15 & 10.41 \\
$n=20$ & Single screen & 20786.00 & 32492.50 & 22695.15 & 10.70 \\
$n=20$ & Adaptive depth-2 screen & \textbf{19535.33} & \textbf{30542.50} & \textbf{21331.24} & 20.67 \\
$n=20$ & \method{} & \textbf{19535.33} & \textbf{30542.50} & \textbf{21331.24} & 77.10 \\
$n=20$ & Screen-gated \method{} & \textbf{19535.33} & \textbf{30542.50} & \textbf{21331.24} & 20.71 \\
\bottomrule
\end{tabular}
\end{table}

\subsection{结构级 AI 与门控}

Depth policy 在 $n=14,16,18$ 项集上训练，并在 held-out $n=20$ 项集上评估。表~\ref{tab:policy} 显示，policy 相对 single screen 获得 42/0/6，平均 score 降低 6.57\%；相对 depth-1 screen 获得 34/0/14，平均 score 降低 2.57\%。相对 all-depth adaptive，policy 平均 score 只高 0.28\%，但平均运行时间降低 34.71\%。这说明结构级 AI 已能学习 screen 深度选择，但尚未在 score 上超过 fixed depth-2。

\begin{table}[t]
\centering
\small
\caption{结构级 depth policy 在 held-out $n=20$ 项集上的比较。}
\label{tab:policy}
\begin{tabular}{lrrrr}
\toprule
比较 & 样本数 & score 胜/负/平 & 平均 $\Delta$ score & 平均 $\Delta$ time \\
\midrule
Policy vs single screen & 48 & 42/0/6 & $-6.57\%$ & $+2224.00\%$ \\
Policy vs depth-1 screen & 48 & 34/0/14 & $-2.57\%$ & $+257.27\%$ \\
Policy vs depth-2 screen & 48 & 0/5/43 & $+0.28\%$ & $-10.85\%$ \\
Policy vs all-depth adaptive & 48 & 0/5/43 & $+0.28\%$ & $-34.71\%$ \\
\bottomrule
\end{tabular}
\end{table}

为去掉这一质量缺口，本文训练了保守 depth-2 skip guard。Static-direct 模式在 held-out $n=20$ test 中没有 false skip，96 个样本全部与 fixed depth-2 screen score 持平，并相对 fixed depth-2 与 all-depth adaptive 分别节省 0.54\% 和 24.74\% 平均时间。Screen-gate 的独立 held-out 验证也支持保守门控：在 $n=19,20$ 共 16 个函数上，screen-gated \method{} 与完整 \method{} 全部 score 持平，平均节省 36.83\% 时间；若直接用 adaptive screen 替代 Resource tail，$n=19$ 子集会出现 4/8 个 score loss。因此，门控贡献不是简单跳过搜索，而是在保守阈值下识别何时可以安全跳过。

\subsection{项集级大规模验证}

表~\ref{tab:scale} 给出 $n=20,22,24,28$ 的项集级 screen-scale 结果。All-depth adaptive screen 相对 single screen 获得 169/0/23，平均 score 降低 6.63\%；但它与 fixed depth-2 在 score 上全部持平，并增加 47.92\% 平均时间。Depth policy 相对 single screen 获得 169/0/23，平均 score 降低 6.56\%；相对 all-depth adaptive 只有 0/6/186，平均 score 增加 0.08\%，但平均时间降低 31.01\%。在最大维度 $n=28$ 上，depth policy 与 all-depth adaptive 在 48 个项集上全部 score 持平，省时 32.15\%。

\begin{table}[H]
\centering
\scriptsize
\caption{$n=20,22,24,28$ 项集级 screen-scale 比较。所有方法行均通过 plan 和 emitted-circuit ANF 符号验证。}
\label{tab:scale}
\begin{tabular}{rllrrr}
\toprule
$n$ & 方法 & baseline & score 胜/负/平 & 平均 $\Delta$ score & 平均 $\Delta$ time \\
\midrule
all & adaptive all-depth & single screen & 169/0/23 & $-6.63\%$ & $+3077.71\%$ \\
all & adaptive all-depth & fixed depth-2 & 0/0/192 & $+0.00\%$ & $+47.92\%$ \\
all & depth policy & single screen & 169/0/23 & $-6.56\%$ & $+2328.66\%$ \\
all & depth policy & all-depth adaptive & 0/6/186 & $+0.08\%$ & $-31.01\%$ \\
28 & depth policy & all-depth adaptive & 0/0/48 & $+0.00\%$ & $-32.15\%$ \\
all & direct depth-2 guard & fixed depth-2 & 0/0/192 & $+0.00\%$ & $-0.14\%$ \\
\bottomrule
\end{tabular}
\end{table}

为检验该结构策略是否只在 $n=20,22,24,28$ 成立，本文进一步使用同一协议扩展到 $n=32,36,40$。表~\ref{tab:scale-extended} 显示，depth policy 相对 single screen 获得 110/0/34，平均 score 降低 5.55\%；相对 all-depth adaptive 在 144 个项集上全部 score 持平，并节省 33.14\% 平均时间。该结果比 $n=20,22,24,28$ 更强地说明，训练于 $n=14,16,18$ 的结构级策略可以外推到更高维项集合。

\begin{table}[H]
\centering
\scriptsize
\caption{$n=32,36,40$ extended screen-scale 比较。所有方法行均通过 plan 和 emitted-circuit ANF 符号验证。}
\label{tab:scale-extended}
\begin{tabular}{rllrrr}
\toprule
$n$ & 方法 & baseline & score 胜/负/平 & 平均 $\Delta$ score & 平均 $\Delta$ time \\
\midrule
all & adaptive all-depth & single screen & 110/0/34 & $-5.55\%$ & $+2680.60\%$ \\
all & adaptive all-depth & fixed depth-2 & 0/0/144 & $+0.00\%$ & $+64.37\%$ \\
all & depth policy & single screen & 110/0/34 & $-5.55\%$ & $+2061.11\%$ \\
all & depth policy & all-depth adaptive & 0/0/144 & $+0.00\%$ & $-33.14\%$ \\
all & direct depth-2 guard & fixed depth-2 & 0/0/144 & $+0.00\%$ & $-0.00\%$ \\
\bottomrule
\end{tabular}
\end{table}

前述策略偏向快速或中等预算。为判断 fixed depth-2 是否已经达到质量上限，本文又在 $n=20,28,40$ 的 72 个项集上运行 depth-3/4 \screen{}。表~\ref{tab:depth-frontier} 显示，depth-3 相对 fixed depth-2 获得 49/0/23，平均 score 降低 1.93\%；depth-4 相对 fixed depth-2 同样获得 49/0/23，平均 score 降低 3.10\%。这说明更深布尔环筛选确实能形成新的质量前沿，但代价是运行时间分别增加 193.71\% 和 682.97\%。因此，depth-3/4 应被写成高预算质量模式，而不是默认快速策略。

\begin{table}[htbp]
\centering
\scriptsize
\caption{$n=20,28,40$ depth-frontier screen-scale 比较。所有方法行均通过 plan 和 emitted-circuit ANF 符号验证。}
\label{tab:depth-frontier}
\begin{tabular}{rllrrr}
\toprule
$n$ & 方法 & baseline & score 胜/负/平 & 平均 $\Delta$ score & 平均 $\Delta$ time \\
\midrule
all & screen depth-3 & fixed depth-2 & 49/0/23 & $-1.93\%$ & $+193.71\%$ \\
all & screen depth-4 & fixed depth-2 & 49/0/23 & $-3.10\%$ & $+682.97\%$ \\
all & screen depth-4 & screen depth-3 & 44/0/28 & $-1.21\%$ & $+127.55\%$ \\
all & all-depth adaptive depth$\leq4$ & fixed depth-2 & 49/0/23 & $-3.10\%$ & $+1129.85\%$ \\
\bottomrule
\end{tabular}
\end{table}

主 scale、extended scale 和 depth-frontier 三组项集级结果合计包含 3000 个方法行。所有方法行均通过两层验证：factor plan 的 \anf{} 展开验证为 3000/3000，emitted X/CNOT/MCT circuit 的 GF(2) 多项式符号模拟验证为 3000/3000，mismatch 总数为 0。这使项集级结果不只是资源估计，但仍应明确它不是完整 \ttable{} simulation。

\section{讨论}

本文最稳妥的贡献表述有四点。第一，方法主线独立于 SSHR。本文不是把 SSHR 改成 AI 版本，而是在 \anf{}/\fprm{} 项集合上做资源加权搜索；SSHR 的作用是 CNOT-oriented baseline。第二，布尔环线性因子提供了明确的结构扩展，允许线性因子与 quotient 共享变量，并在 $n=18,20$ 上稳定改善 single screen。第三，结构级 AI 已有正向证据：depth policy 学会了 screen 深度选择，保守 guard 消除了相对 fixed depth-2 的 score 缺口，并减少 all-depth adaptive overhead。第四，项集级 scale 测试将证据推进到 $n=20,22,24,28,32,36,40$，depth-frontier 进一步显示 depth-3/4 可以打破 fixed depth-2 的质量上限，并给出 emitted-circuit 级的符号等价验证。

当前结果还不能支撑更强的结论。其一，$n=18$ 和 held-out $n=19$ 上完整 \method{} 仍能超过 adaptive screen，说明中等高维仍需要更深搜索；其二，$n=20$ 切片中完整 \method{} 与 adaptive screen 资源持平，说明该切片的收益主要来自结构筛选而非 MCTS tail；其三，高维 learned prior 的独立质量收益仍弱，当前 AI 主体更多体现为结构选择和安全门控；其四，$n=22$ 到 $n=40$ 的项集级结果虽有 emitted-circuit ANF 符号等价验证，但不包含完整 \ttable{} 枚举。

下一步应把论文推进到三个可量化目标。第一，把 shallow-staged guard 的安全跳过能力蒸馏回 direct guard，使 held-out $n=20$ 上继续保持不劣于 fixed depth-2 的 score，同时把 direct 模式相对 fixed depth-2 的时间优势提高到至少 5\%。第二，把 screen-gate 从单一边界规则扩展为多特征分支选择器，联合预测 screen 深度、是否进入 \fprm{} polarity search、是否运行完整 \method{} tail。第三，补充更多外部 reversible-toolchain 或 ROS/mockturtle 风格比较，使论文从内部资源表推进到更有说服力的工具链对比。

\section{结论}

本文提出了面向资源约束量子布尔函数综合的 \method{} 方法。它以 \anf{}/\fprm{} 项集合为状态，以逻辑资源 score 为目标，用布尔环线性因子、神经动作先验、\mcts{}、baseline guard、depth policy 和 screen-gate 共同生成候选线路。实验表明，该方法在 $n\leq6$ 传统函数上显著降低 T-count 与加权 score，在 $n=18,20$ 高维函数级切片上验证了 adaptive screen 与 Resource tail 的互补作用，在 $n=20$ 到 $n=40$ 项集级 scale 测试中展示了结构级 AI 的可扩展性，并通过 depth-frontier 实验证明更深 Boolean-ring screen 可在高预算下继续降低 score。当前稿件已经形成独立于 SSHR 的论文主线，但最终投稿还需要继续强化高维神经策略的质量收益、完整 \ttable{} 可验证切片和外部工具链比较。

\begin{thebibliography}{99}

\bibitem{meuli2019multiplicative}
G.~Meuli, M.~Soeken, E.~Campbell, M.~Roetteler, and G.~De Micheli.
\newblock The role of multiplicative complexity in compiling low T-count oracle circuits.
\newblock In \emph{Proceedings of ICCAD}, 2019.

\bibitem{meuli2022xag}
G.~Meuli, M.~Soeken, and G.~De Micheli.
\newblock Xor-And-Inverter Graphs for quantum compilation.
\newblock \emph{npj Quantum Information}, 8:7, 2022.

\bibitem{meuli2020ros}
G.~Meuli, M.~Soeken, M.~Roetteler, and G.~De Micheli.
\newblock ROS: Resource-constrained oracle synthesis for quantum computers.
\newblock \emph{Electronic Proceedings in Theoretical Computer Science}, 318:119--130, 2020.

\bibitem{wille2009bdd}
R.~Wille and R.~Drechsler.
\newblock BDD-based synthesis of reversible logic for large functions.
\newblock In \emph{Proceedings of DAC}, pp.~270--275, 2009.

\bibitem{yu2025backend}
M.~Yu, A.~Tempia Calvino, M.~Soeken, and G.~De Micheli.
\newblock Back-end-aware fault-tolerant quantum oracle synthesis.
\newblock In \emph{Proceedings of ASP-DAC}, pp.~930--937, 2025.

\bibitem{zheng2025sshr}
Q.~Zheng, Y.~Xu, J.~Zhang, Z.~Su, and S.~Zheng.
\newblock CNOT oriented synthesis for small-scale Boolean functions using spatial structures of parallelotopes.
\newblock In \emph{Proceedings of ICCAD}, 2025.

\bibitem{tsaras2024shortcircuit}
D.~Tsaras, A.~Grosnit, L.~Chen, Z.~Xie, H.~Bou-Ammar, and M.~Yuan.
\newblock ShortCircuit: AlphaZero-driven circuit design.
\newblock \emph{arXiv:2408.09858}, 2024.

\bibitem{rietsch2024unitary}
S.~Rietsch, A.~Y. Dubey, C.~Ufrecht, M.~Periyasamy, A.~Plinge, C.~Mutschler, and D.~D. Scherer.
\newblock Unitary synthesis of Clifford+T circuits with reinforcement learning.
\newblock In \emph{IEEE International Conference on Quantum Computing and Engineering}, pp.~824--835, 2024.

\bibitem{riu2025rlzx}
J.~Riu, J.~Nogu{\'e}, G.~Vilaplana, A.~Garcia-Saez, and M.~P. Estarellas.
\newblock Reinforcement learning based quantum circuit optimization via ZX-calculus.
\newblock \emph{Quantum}, 9:1758, 2025.

\bibitem{machiya2026monteq}
M.~Machiya, M.~Menickelly, P.~Hovland, and J.~Liu.
\newblock MonteQ: A Monte Carlo tree search based quantum circuit synthesis framework.
\newblock \emph{arXiv:2604.19029}, 2026.

\end{thebibliography}

\end{document}
